\documentclass[11pt]{article}
\usepackage{fontspec,amsmath,amssymb,amsthm,hyperref}
\usepackage[margin=1in]{geometry}
\setlength{\emergencystretch}{3em}
\title{Fourth moments of actual time integrals}
\author{LangevinLean mathematical exposition}
\date{}
\begin{document}\maketitle
\noindent Source module: \texttt{KineticLangevin/LpMomentIntegral.lean}. Compile this proof document with XeLaTeX or LuaLaTeX. The accompanying HTML pairs the mathematical argument with the exact Lean statements and records the scope of the formalization.


Kinetic positions are time integrals of velocities in L\ensuremath{{}^{2}}. A fourth-moment estimate must apply to this existing L\ensuremath{{}^{2}} integral, rather than to a separately postulated L4 process. The proof uses closed convex L4 balls inside L\ensuremath{{}^{2}} and the geometric fact that averages remain in a closed convex set.


\section{View an L4 ball inside L\ensuremath{{}^{2}}}


The set lpFourthBall(R) contains actual L\ensuremath{{}^{2}} equivalence classes whose extended fourth norm is at most R. If a sequence in this set converges in L\ensuremath{{}^{2}}, it converges in measure. Lower semicontinuity of the extended L4 norm under convergence in measure keeps its limit in the same ball. In the metric L\ensuremath{{}^{2}} space, this sequential argument proves closedness.



\noindent Lean declarations: \texttt{lpFourthBall}, \texttt{isClosed\_lpFourthBall}.

\section{Verify convexity without extra moment assumptions}


For nonnegative a,b with a+b=1, the actual representative of aX+bY equals aX(ω)+bY(ω) almost everywhere. The L4 triangle inequality and homogeneity give ‖aX+bY‖₄\ensuremath{\le}a‖X‖₄+b‖Y‖₄\ensuremath{\le}R. Extended norms allow this argument even before finite fourth moments are separately established.



\noindent Lean declarations: \texttt{convex\_lpFourthBall}.

\section{Average a continuous L\ensuremath{{}^{2}} path}


If every velocity v(t) on [0,T] has fourth norm at most R, its time average lies in the closed convex L4 ball. The average is T⁻¹ times the actual Bochner integral in L\ensuremath{{}^{2}}. Homogeneity then gives ‖∫₀ᵀv(t)dt‖₄\ensuremath{\le}TR. The zero-duration case is handled separately; for positive duration Lean checks the interval measure is finite and nonzero before applying the convex averaging theorem. This directly supplies displacement fourth moments from Gaussian velocity marginals.



\noindent Lean declarations: \texttt{eLpNorm\_four\_integral\_le}.

\end{document}
